\chapter{A derivada} \label{der:chapter}

%%%%%%%%%%%%%%%%%%%%%%%%%%%%%%%%%%%%%%%%%%%%%%%%%%%%%%%%%%%%%%%%%%%%%%%%%%%%%%

\section{A derivada}
\label{sec:der}

%mbxINTROSUBSECTION

\sectionnotes{1 aula}

A ideia de derivada é a seguinte.
Se o gráfico de uma função se parece localmente com uma reta,
podemos falar da inclinação dessa reta.
A inclinação indica a taxa de variação do valor da função naquele ponto.
Naturalmente, estamos deixando de lado as funções que têm bicos ou
descontinuidades.
Sejamos precisos.

\subsection{Definição e propriedades básicas}

\begin{defn}
Seja $I$ um intervalo, seja
$f \colon I \to \R$ uma função e seja $c \in I$. Se o limite
\begin{equation*}
L \coloneqq \lim_{x \to c} \frac{f(x)-f(c)}{x-c} 
\end{equation*}
existe, dizemos que $f$ é
\emph{\myindex{diferenciável}}\index{função!diferenciável} em
$c$, chamamos $L$ de \emph{\myindex{derivada}} de $f$ em $c$
e escrevemos $f'(c) \coloneqq L$.\glsadd{not:derivative}

\medskip

Se $f$ é diferenciável em todo $c \in I$, dizemos simplesmente que
$f$ é \emph{diferenciável}; assim, obtemos uma função
$f' \colon I \to \R$.
Às vezes, a derivada é denotada por $\frac{df}{dx}$ ou
$\frac{d}{dx}\bigl( f(x) \bigr)$.

\medskip

A expressão $\frac{f(x)-f(c)}{x-c}$ é chamada de
\emph{\myindex{quociente de diferenças}}.
\end{defn}

A interpretação gráfica da derivada está ilustrada na
\figureref{derivfig}. O gráfico à esquerda mostra a reta que passa por
$\bigl(c,f(c)\bigr)$ e $\bigl(x,f(x)\bigr)$, com inclinação
$\frac{f(x)-f(c)}{x-c}$; essa é a
\emph{\myindex{reta secante}}. Ao tomarmos o limite quando $x$ tende a $c$,
obtemos o gráfico à direita, em que vemos que a derivada da função
no ponto $c$ é a inclinação da reta tangente ao gráfico de $f$
no ponto $\bigl(c,f(c)\bigr)$.

\begin{myfigureht}
\myincludepdft{deriv_derivd}{%
Dois diagramas do gráfico de uma função f. À esquerda,
os pontos c e x estão marcados no eixo horizontal.
Os pontos correspondentes no gráfico de f também estão marcados,
e uma reta secante passa por eles. Sua inclinação está indicada por
f de x menos f de c, tudo dividido por x menos c.
À direita, a mesma função está representada, mas apenas o ponto c
está marcado. Uma reta tangente ao gráfico de f passa pelo ponto
correspondente e está identificada com inclinação igual a f linha de c.}
\caption{Interpretação gráfica da derivada.\label{derivfig}}
\end{myfigureht}

Permitimos que $I$ seja um intervalo fechado e que
$c$ seja uma extremidade de $I$. Alguns livros de cálculo não permitem que
$c$ seja uma extremidade do intervalo, mas toda a teoria continua válida
quando a permitimos, e isso simplifica nosso trabalho.

\begin{example}
Seja $f(x) \coloneqq x^2$ definida em toda a reta real. Seja $c \in \R$
arbitrário. Vemos que, se $x \neq c$,
\begin{equation*}
\frac{x^2-c^2}{x-c} =
\frac{(x+c)(x-c)}{x-c} =
(x+c) .
\end{equation*}
Portanto,
\begin{equation*}
f'(c) =
\lim_{x\to c} \frac{x^2-c^2}{x-c} =
\lim_{x\to c} (x+c) = 2c.
\end{equation*}
\end{example}

\begin{example}
Seja $f(x) \coloneqq ax + b$, em que $a, b \in \R$.
Seja $c \in \R$ arbitrário.
Para $x \neq c$,
\begin{equation*}
\frac{f(x)-f(c)}{x-c} =
\frac{a(x-c)}{x-c} = a .
\end{equation*}
Portanto,
\begin{equation*}
f'(c) =
\lim_{x\to c}
\frac{f(x)-f(c)}{x-c} =
\lim_{x\to c}
a = a.
\end{equation*}
De fato, toda função diferenciável se comporta \myquote{infinitesimalmente}
como a função afim $ax + b$. É possível conjecturar muitos resultados
e fórmulas para derivadas calculando-as primeiro para funções afins.
\end{example}

\begin{example}
A função $f(x) \coloneqq \sqrt{x}$ é diferenciável para $x > 0$. Para ver
isso, fixe $c > 0$ e suponha que $x \neq c$ e $x > 0$. Calcule
\begin{equation*}
\frac{\sqrt{x}-\sqrt{c}}{x-c}
=
\frac{\sqrt{x}-\sqrt{c}}{(\sqrt{x}-\sqrt{c})(\sqrt{x}+\sqrt{c})}
=
\frac{1}{\sqrt{x}+\sqrt{c}} .
\end{equation*}
Portanto,
\begin{equation*}
f'(c) =
\lim_{x\to c}
\frac{\sqrt{x}-\sqrt{c}}{x-c}
=
\lim_{x\to c}
\frac{1}{\sqrt{x}+\sqrt{c}}
=
\frac{1}{2\sqrt{c}} .
\end{equation*}
\end{example}

\begin{example}
A função $f(x) \coloneqq \sabs{x}$ não é diferenciável
na origem. Quando $x > 0$,
\begin{equation*}
\frac{\sabs{x}-\sabs{0}}{x-0} =
\frac{x-0}{x-0} = 1 .
\end{equation*}
Quando $x < 0$,
\begin{equation*}
\frac{\sabs{x}-\sabs{0}}{x-0} =
\frac{-x-0}{x-0} = -1 .
\end{equation*}
\end{example}

Um famoso exemplo de Weierstrass mostra que existe uma função contínua
que não é diferenciável em \emph{nenhum} ponto. A construção dessa função
está além do escopo deste capítulo. Por outro lado, toda função diferenciável
é contínua.

\begin{prop}
Se $f \colon I \to \R$ é diferenciável em $c \in I$,
então é contínua em $c$.
\end{prop}

\begin{proof}
Sabemos que os limites
\begin{equation*}
\lim_{x\to c}\frac{f(x)-f(c)}{x-c} = f'(c)
\qquad
\text{e}
\qquad
\lim_{x\to c}(x-c) = 0
\end{equation*}
existem. Além disso,
\begin{equation*}
f(x)-f(c) = 
\left( \frac{f(x)-f(c)}{x-c} \right) (x-c) .
\end{equation*}
Portanto, o limite de $f(x)-f(c)$ existe e
\begin{equation*}
\lim_{x\to c} \bigl( f(x)-f(c) \bigr) =
\left(\lim_{x\to c} \frac{f(x)-f(c)}{x-c} \right)
\left(\lim_{x\to c} (x-c) \right) =
f'(c) \cdot 0  = 0.
\end{equation*}
Logo, $\lim\limits_{x\to c} f(x) = f(c)$, e $f$ é contínua em $c$.
\end{proof}

Uma propriedade importante da derivada é a linearidade. A derivada aproxima
uma função por uma reta. A inclinação da reta que passa por dois pontos
varia linearmente quando as coordenadas $y$ variam linearmente.
Ao tomarmos o limite, faz sentido que a derivada seja linear.

\begin{prop}[Linearidade]
\index{linearidade da derivada}
Seja $I$ um intervalo, sejam
$f \colon I \to \R$ e $g \colon I \to \R$ diferenciáveis em $c \in I$,
e seja $\alpha \in \R$.
\begin{enumerate}[(i)]
\item
Defina $h \colon I \to \R$ por $h(x) \coloneqq \alpha f(x)$. Então
$h$ é diferenciável em $c$ e
$h'(c) = \alpha f'(c)$.
\item
Defina $h \colon I \to \R$ por $h(x) \coloneqq f(x) + g(x)$. Então
$h$ é diferenciável em $c$ e
$h'(c) =  f'(c) + g'(c)$.
\end{enumerate}
\end{prop}

\begin{proof}
Primeiro, seja $h(x) \coloneqq \alpha f(x)$.
Para $x \in I$, $x \neq c$,
\begin{equation*}
\frac{h(x)-h(c)}{x-c} =
\frac{\alpha f(x) - \alpha f(c)}{x-c}
=
\alpha \frac{f(x) - f(c)}{x-c} .
\end{equation*}
O limite quando $x$ tende a $c$ existe no lado direito
por \corref{falg:cor}. Assim,
\begin{equation*}
\lim_{x\to c}\frac{h(x)-h(c)}{x-c} =
\alpha \lim_{x\to c} \frac{f(x) - f(c)}{x-c} .
\end{equation*}
Portanto, $h$ é diferenciável em $c$,
e a derivada é calculada como indicado.

Em seguida, defina $h(x) \coloneqq f(x)+g(x)$.
Para $x \in I$, $x \neq c$, temos
\begin{equation*}
\frac{h(x)-h(c)}{x-c} =
\frac{\bigl(f(x) + g(x)\bigr) - \bigl(f(c) + g(c)\bigr)}{x-c}
=
\frac{f(x) - f(c)}{x-c}
+
\frac{g(x) - g(c)}{x-c} .
\end{equation*}
O limite quando $x$ tende a $c$ existe no lado direito
por \corref{falg:cor}. Assim,
\begin{equation*}
\lim_{x\to c}\frac{h(x)-h(c)}{x-c} =
\lim_{x\to c} \frac{f(x) - f(c)}{x-c}
+
\lim_{x\to c}\frac{g(x) - g(c)}{x-c} .
\end{equation*}
Portanto, $h$ é diferenciável em $c$,
e a derivada é calculada como indicado.
\end{proof}

A derivada do produto de duas funções não é o produto das derivadas.
Em vez disso, obtemos o que se chama \emph{regra do produto}
ou \emph{\myindex{regra de Leibniz}}%
\footnote{Em homenagem ao matemático alemão
\href{https://en.wikipedia.org/wiki/Leibniz}{Gottfried Wilhelm Leibniz}
(1646--1716).}.

\begin{prop}[Regra do produto]\index{regra do produto}
Seja $I$ um intervalo, sejam
$f \colon I \to \R$ e $g \colon I \to \R$ funções diferenciáveis em $c$.
Se $h \colon I \to \R$ é definida por
\begin{equation*}
h(x) \coloneqq f(x) g(x) ,
\end{equation*}
então $h$ é diferenciável em $c$ e
\begin{equation*}
h'(c) = f(c) g'(c) + f'(c) g(c) .
\end{equation*}
\end{prop}

A demonstração da regra do produto fica como exercício. A identidade
fundamental para a demonstração é
$f(x) g(x) - f(c) g(c) =
f(x)\bigl( g(x) - g(c) \bigr)
+ \bigl( f(x) - f(c) \bigr) g(c)$,
como ilustra a \figureref{figprodrule}.
\begin{myfigureht}
\myincludepdft{figprodrule}{%
Diagrama de uma região retangular no plano, dividida
em três retângulos. O lado horizontal vai de 0 a
f de x, e o lado vertical vai de 0 a g de x.
A maior região branca é um retângulo que se estende horizontalmente
de 0 até f de c, um pouco menor que f de x,
e verticalmente de 0 até g de c, um pouco menor que
g de x. Essa região branca está identificada como f de c vezes g de c.
Acima dela, há um retângulo fino, levemente sombreado, que se estende
horizontalmente de 0 a f de x e verticalmente
de g de c a g de x. Ele está identificado
como f de x vezes a quantidade g de x menos g de c.
Por fim, há um retângulo fino, mais escuro, que se estende
horizontalmente de f de c a f de x e verticalmente
de 0 a g de c. Ele está identificado como f de x menos f de c,
tudo isso vezes g de c.}
\caption{Ideia da regra do produto. A área do retângulo inteiro
$f(x)g(x)$ difere da área do retângulo não sombreado $f(c)g(c)$
pela área dos dois retângulos sombreados:
$f(x)\bigl( g(x) - g(c) \bigr)$ e
$\bigl( f(x) - f(c) \bigr) g(c)$.
Em outras palavras, aproximadamente, $\Delta (f \cdot g)
= f \cdot \Delta g + \Delta f \cdot g$.\label{figprodrule}}
\end{myfigureht}



\begin{prop}[Regra do quociente]\index{regra do quociente}
Seja $I$ um intervalo, sejam
$f \colon I \to \R$ e $g \colon I \to \R$ diferenciáveis em $c$,
e suponha que $g(x) \neq 0$ para todo $x \in I$.
Se $h \colon I \to \R$ é definida por
\begin{equation*}
h(x) \coloneqq \frac{f(x)}{g(x)},
\end{equation*}
então $h$ é diferenciável em $c$ e
\begin{equation*}
h'(c) = \frac{f'(c) g(c) - f(c) g'(c)}{{\bigl(g(c)\bigr)}^2} .
\end{equation*}
\end{prop}

Novamente, a demonstração fica como exercício.

\subsection{Regra da cadeia}

Funções mais complicadas são frequentemente obtidas por composição,
e sua derivada é calculada pela regra da cadeia. Essa regra também nos diz
como a derivada muda quando fazemos uma mudança de variável.

\begin{prop}[Regra da cadeia]
\index{regra da cadeia}
Sejam $I_1, I_2$ intervalos, seja
$g \colon I_1 \to I_2$ diferenciável em $c \in I_1$
e seja
$f \colon I_2 \to \R$ diferenciável em $g(c)$.
Se $h \colon I_1 \to \R$ é definida por
\begin{equation*}
h(x) \coloneqq (f \circ g) (x) = f\bigl(g(x)\bigr) ,
\end{equation*}
então $h$ é diferenciável em $c$ e
\begin{equation*}
h'(c) = f'\bigl(g(c)\bigr)g'(c) .
\end{equation*}
\end{prop}

\begin{proof}
Seja $d \coloneqq g(c)$. Defina
$u \colon I_2 \to \R$ e $v \colon I_1 \to \R$ por
\begin{equation*}
u(y) \coloneqq
\begin{cases}
 \frac{f(y) - f(d)}{y-d}  & \text{se } y \neq d, \\
 f'(d)                    & \text{se } y = d,
\end{cases}
\qquad
v(x) \coloneqq
\begin{cases}
 \frac{g(x) - g(c)}{x-c} & \text{se } x \neq c, \\
 g'(c)                   & \text{se } x = c.
\end{cases}
\end{equation*}
Como $f$ é diferenciável em $d = g(c)$, concluímos que
$u$ é contínua em $d$. De modo semelhante, $v$ é contínua em $c$.
Para quaisquer $x$ e $y$,
\begin{equation*}
f(y)-f(d) = u(y) (y-d)
\qquad \text{e} \qquad
g(x)-g(c) = v(x) (x-c) .
\end{equation*}
Substituindo, obtemos
\begin{equation*}
h(x)-h(c)
=
f\bigl(g(x)\bigr)-f\bigl(g(c)\bigr)
=
u\bigl( g(x) \bigr) \bigl(g(x)-g(c)\bigr)
=
u\bigl( g(x) \bigr) \bigl(v(x) (x-c)\bigr) .
\end{equation*}
Portanto, se $x \neq c$,
\begin{equation} \label{eq:chainruleeq}
\frac{h(x)-h(c)}{x-c}
=
u\bigl( g(x) \bigr) v(x) .
\end{equation}
Pela continuidade de $u$ e $v$ em $d$ e $c$, respectivamente, temos
$\lim_{y \to d} u(y)
= f'(d) = f'\bigl(g(c)\bigr)$ e
$\lim_{x \to c} v(x) = g'(c)$.
A função $g$ é contínua em $c$ e, portanto, $\lim_{x \to c} g(x) = g(c)$.
Assim, o limite do lado direito de \eqref{eq:chainruleeq}
quando $x$ tende a $c$ existe e é igual a $f'\bigl(g(c)\bigr) g'(c)$. Logo, $h$
é diferenciável em $c$ e $h'(c) = f'\bigl(g(c)\bigr)g'(c)$.
\end{proof}

\subsection{Exercícios}

\begin{exercise}
Demonstre a regra do produto.
Dica: demonstre e use
$f(x) g(x) - f(c) g(c) = f(x)\bigl( g(x) - g(c) \bigr) + \bigl( f(x) -
f(c) \bigr) g(c)$.
\end{exercise}

\begin{exercise}
Demonstre a regra do quociente. Dica: você pode fazer isso diretamente, mas
talvez seja mais fácil encontrar a derivada de $\nicefrac{1}{x}$ e, em seguida,
usar a regra da cadeia e a regra do produto.
\end{exercise}

\begin{exercise} \label{exercise:diffofxn}
Para $n \in \Z$, demonstre que $x^n$ é diferenciável e encontre sua derivada,
exceto, naturalmente, quando $n < 0$ e $x=0$.
Dica: use a regra do produto.
\end{exercise}

\begin{exercise}
Demonstre que um polinômio é diferenciável e encontre sua derivada.
Dica: use o exercício anterior.
\end{exercise}

\begin{exercise}
Defina $f \colon \R \to \R$ por
\begin{equation*}
f(x) \coloneqq
\begin{cases}
x^2 & \text{se } x \in \Q,\\
0 & \text{caso contrário.}
\end{cases}
\end{equation*}
Demonstre que $f$ é diferenciável em $0$, mas descontínua em todos os pontos,
exceto em $0$.
\end{exercise}

\begin{exercise}
Suponha que vale a desigualdade $\babs{x-\sin(x)} \leq x^2$. Demonstre que
$\sin$ é diferenciável em $0$ e encontre a derivada em $0$.
\end{exercise}

\begin{exercise}
Usando o exercício anterior, demonstre que $\sin$ é diferenciável para todo $x$
e que sua derivada é $\cos(x)$. Dica: use a identidade trigonométrica
de soma em produto, como fizemos anteriormente.
\end{exercise}

\begin{exercise}
Seja $f \colon I \to \R$ diferenciável. Para $n \in \Z$, seja $f^n$
a função definida por $f^n(x) \coloneqq {\bigl( f(x) \bigr)}^n$. Se
$n < 0$, suponha que $f(x) \neq 0$ para todo $x \in I$. Demonstre que
$(f^n)'(x) = n {\bigl(f(x) \bigr)}^{n-1} f'(x)$.
\end{exercise}

\begin{exercise}
Suponha que $f \colon \R \to \R$ seja uma função diferenciável
e Lipschitz contínua.
Demonstre que $f'$ é uma função limitada.
\end{exercise}

\begin{exercise} \label{exercise:inversederformula}
Sejam $I_1, I_2$ intervalos.
Seja $f \colon I_1 \to I_2$ uma função bijetiva e seja $g \colon I_2 \to I_1$
sua inversa. Suponha que $f$ seja diferenciável em $c \in I_1$,
que $f'(c) \neq 0$ e que $g$ seja diferenciável em $f(c)$. Use a regra da cadeia
para encontrar uma fórmula para $g'\bigl(f(c)\bigr)$ (em termos de $f'(c)$).
\end{exercise}

\begin{exercise} \label{exercise:bndmuldiff}
Suponha que $f \colon I \to \R$ seja limitada, que $g \colon I \to
\R$ seja diferenciável em $c \in I$ e que $g(c) = g'(c) = 0$. Mostre
que $h(x) \coloneqq f(x) g(x)$ é diferenciável em $c$. Dica: não é possível
aplicar a regra do produto.
\end{exercise}

\begin{exercise} \label{exercise:diffsqueeze}
Suponha que $f \colon I \to \R$,
$g \colon I \to \R$ e $h \colon I \to \R$ sejam funções. Suponha que
$c \in I$ satisfaça $f(c) = g(c) = h(c)$, que $g$ e $h$ sejam diferenciáveis
em $c$ e que $g'(c) = h'(c)$. Além disso, suponha que $h(x) \leq f(x) \leq g(x)$
para todo $x \in I$. Demonstre que $f$ é diferenciável em $c$ e que
$f'(c) = g'(c) = h'(c)$.
\end{exercise}

\begin{exercise}
\leavevmode
\begin{enumerate}[a)]
\item
Suponha que $f \colon (-1,1) \to \R$ seja uma função tal que $f(x) = x h(x)$,
em que $h \colon (-1,1) \to \R$ é limitada.
Mostre que $g(x) \coloneqq {\bigl( f(x) \bigr)}^2$ é diferenciável na origem
e que $g'(0) = 0$.
\item
Encontre um exemplo de função contínua $f \colon (-1,1) \to \R$ com $f(0) = 0$,
mas tal que $g(x) \coloneqq {\bigl( f(x) \bigr)}^2$ não seja diferenciável na origem.
\end{enumerate}
\end{exercise}

\begin{exercise}
Suponha que $f \colon I \to \R$ seja diferenciável em $c \in I$.
Demonstre que existem números $a$ e $b$ tais que, para todo $\epsilon > 0$,
existe $\delta > 0$ com a propriedade de que
$\babs{a+b(x-c) - f(x)} \leq \epsilon \sabs{x-c}$ sempre que $x \in I$ e
$\sabs{x-c} < \delta$.
Em outras palavras, mostre que existe uma função $g \colon I \to \R$
tal que $\lim_{x\to c} g(x) = 0$ e
$\babs{a+b(x-c) - f(x)} = g(x) \sabs{x-c}$.
\end{exercise}

\begin{exercise} \label{exercise:simpleLHopital}
Demonstre a seguinte versão simples da \myindex{regra de L'H\^opital}\index{regra de L'Hospital}.
Suponha que
$f \colon (a,b) \to \R$ e $g \colon (a,b) \to \R$ sejam funções diferenciáveis
cujas derivadas $f'$ e $g'$ são contínuas.
Suponha que, em $c \in (a,b)$, $f(c) = 0$ e $g(c)=0$,
que $g'(x) \neq 0$ para todo $x \in (a,b)$ e que
$g(x) \neq 0$ sempre que $x \neq c$. Observe
que existe o limite de $\nicefrac{f'(x)}{g'(x)}$ quando $x$ tende a $c$. Mostre que
\begin{equation*}
\lim_{x \to c} \frac{f(x)}{g(x)} = 
\lim_{x \to c} \frac{f'(x)}{g'(x)} .
\end{equation*}
\end{exercise}

\begin{exercise}
Suponha que $f \colon (a,b) \to \R$ seja diferenciável em $c \in (a,b)$,
$f(c)=0$ e $f'(c) > 0$.
Demonstre que existe $\delta > 0$ tal que $f(x) < 0$
sempre que $c-\delta < x < c$ e
$f(x) > 0$
sempre que $c < x < c+\delta$.
\end{exercise}

%%%%%%%%%%%%%%%%%%%%%%%%%%%%%%%%%%%%%%%%%%%%%%%%%%%%%%%%%%%%%%%%%%%%%%%%%%%%%%

\sectionnewpage
\section{Teorema do valor médio}
\label{sec:mvt}

%mbxINTROSUBSECTION

\sectionnotes{2 aulas (algumas aplicações podem ser omitidas)}

\subsection{Mínimos e máximos relativos}

Já falamos sobre máximos e mínimos absolutos. Eles são os picos mais altos
e os vales mais profundos de toda a cadeia de montanhas. E quanto aos picos
de montanhas individuais e aos fundos de vales específicos?
A derivada, por ser um conceito local, é como caminhar em meio à neblina:
ela não pode dizer se você está no pico mais alto, mas pode dizer se está
no topo de algum pico.

\begin{defn}
Seja $S \subset \R$ um conjunto e
seja $f \colon S \to \R$ uma função. Dizemos que $f$ tem
um \emph{\myindex{máximo relativo}}\index{máximo!relativo}
em $c \in S$ se existe $\delta>0$
tal que, para todo $x \in S$ com $\sabs{x-c} < \delta$,
temos $f(x) \leq f(c)$.
A definição de
\emph{\myindex{mínimo relativo}}\index{mínimo!relativo}
é análoga.
\end{defn}

\begin{lemma}\label{relminmax:lemma}
Suponha que $f \colon (a,b) \to \R$ seja diferenciável em $c \in (a,b)$
e que $f$ tenha
um mínimo relativo ou um máximo relativo em $c$. Então
$f'(c) = 0$.
\end{lemma}

\begin{proof}
Suponha que $c$ seja um
máximo relativo de $f$. Isto é, existe $\delta > 0$ tal
que, para todo $x \in (a,b)$ com
$\sabs{x-c} < \delta$, temos $f(x)-f(c) \leq 0$.
Considere o quociente de diferenças. Se $c < x < c+\delta$, então
\begin{equation*}
\frac{f(x)-f(c)}{x-c} \leq 0 ,
\end{equation*}
e, se $c-\delta < y < c$, então
\begin{equation*}
\frac{f(y)-f(c)}{y-c} \geq 0 .
\end{equation*}
Veja na \figureref{fig:critpt} uma ilustração.
\begin{myfigureht}
\myincludegraphics{critpt}{%
Gráfico de uma função f.
Três pontos estão marcados no eixo horizontal, da esquerda
para a direita: y, c e x. A função tem um máximo
no ponto c. Os pontos correspondentes a x, c e y também estão
marcados no gráfico.
Uma reta secante passa pelos pontos correspondentes
a y e c, subindo da esquerda para a direita, e está identificada
com inclinação igual a f de y menos f de c, tudo isso
dividido por y menos c, maior ou igual a 0.
Outra reta secante passa pelos pontos correspondentes
a c e x, descendo da esquerda para a direita, e está identificada
com inclinação igual a f de x menos f de c, tudo isso
dividido por x menos c, menor ou igual a 0.}
\caption{Inclinações de retas secantes em um máximo relativo.\label{fig:critpt}}
\end{myfigureht}

Como $a < c < b$, existem
sequências $\{ x_n \}_{n=1}^\infty$ e
$\{ y_n \}_{n=1}^\infty$ em $(a,b)$
tais que $c < x_n < c+\delta$ e
$c-\delta < y_n < c$ para todo $n \in \N$, e tais que
$\lim_{n\to\infty} x_n = \lim_{n\to\infty} y_n = c$.
Como $f$
é diferenciável em $c$,
\begin{equation*}
0 \geq \lim_{n\to\infty} \frac{f(x_n)-f(c)}{x_n-c} 
=
f'(c)
=
\lim_{n\to\infty} \frac{f(y_n)-f(c)}{y_n-c} \geq 0.
\end{equation*}
Concluímos a demonstração para um máximo.
Para um mínimo, considere
a função $-f$.
\end{proof}

Para uma função diferenciável, chama-se
\emph{\myindex{ponto crítico}} um ponto em que
$f'(c) = 0$. Quando $f$ não é diferenciável em alguns pontos,
também é comum dizer que $c$ é um ponto crítico
se $f'(c)$ não existe.
O teorema afirma que um mínimo ou máximo relativo em um ponto interior
de um intervalo deve ser um ponto crítico.
Como você se lembra do cálculo, para encontrar os mínimos e máximos de uma
função, determinamos todos os pontos críticos e também as extremidades
do intervalo; em seguida, basta verificar em qual desses pontos
a função assume seu maior ou menor valor.

\subsection{Teorema de Rolle}

Suponha que uma função tenha o mesmo valor nas duas extremidades de um intervalo.
Intuitivamente, ela deveria atingir um mínimo ou um máximo no interior do
intervalo.
Nesse mínimo ou máximo, a derivada deveria ser nula.
Veja na \figureref{rollefig} a ideia geométrica.
Esse é o conteúdo do teorema de Rolle%
\footnote{Em homenagem ao matemático francês
\href{https://en.wikipedia.org/wiki/Michel_Rolle}{Michel Rolle}
(1652--1719).}.

\begin{myfigureht}
\myincludegraphics{rollefig}{%
Gráfico de uma função no intervalo de a a b.
A função vale zero em a e em b. Ela cresce até
um máximo, depois decresce até um mínimo e volta a crescer. O máximo da função ocorre
em um ponto marcado como c. Uma reta tangente horizontal está traçada no
ponto correspondente do gráfico. Note que ela é paralela ao
eixo horizontal.}
\caption{Ponto em que a reta tangente é horizontal, isto é, $f'(c) =
0$.\label{rollefig}}
\end{myfigureht}

\begin{thm}[Rolle] \label{thm:rolle}
\index{teorema de Rolle}
Seja $f \colon [a,b] \to \R$ uma função contínua
e diferenciável em $(a,b)$ tal que $f(a) = f(b)$.
Então existe $c \in (a,b)$ tal que $f'(c) = 0$.
\end{thm}

\begin{proof}
Como $f$ é contínua em $[a,b]$, ela atinge um mínimo absoluto e um
máximo absoluto em $[a,b]$. Queremos aplicar o
lema \lemmaref{relminmax:lemma}; portanto,
precisamos encontrar algum $c \in (a,b)$ em que $f$ atinja um mínimo ou
um máximo.
Escreva $K \coloneqq f(a) = f(b)$.
Se existe $x$ tal que $f(x) > K$, então o máximo absoluto é maior que $K$
e, portanto, ocorre em algum $c \in (a,b)$; logo, $f'(c) = 0$.
Por outro lado, se existe $x$ tal que $f(x) < K$, então o mínimo absoluto
ocorre em algum $c \in (a,b)$ e, assim, $f'(c) = 0$. Se não existe $x$ tal que
$f(x) > K$ ou
$f(x) < K$, então $f(x) = K$ para todo $x$ e, portanto,
$f'(x) = 0$ para todo $x \in [a,b]$. Assim, qualquer $c \in (a,b)$ serve.
\end{proof}

É absolutamente necessário que a derivada exista para todo $x
\in (a,b)$. Considere a função $f(x) \coloneqq \sabs{x}$ em $[-1,1]$.
É claro que $f(-1) = f(1)$, mas não existe ponto $c$ em que $f'(c) = 0$.

\subsection{Teorema do valor médio}

Estendemos o \hyperref[thm:rolle]{teorema de Rolle}
a funções que assumem valores diferentes
nas extremidades.

\begin{thm}[Teorema do valor médio] \label{thm:mvt}
\index{teorema do valor médio}
Seja $f \colon [a,b] \to \R$ uma função contínua em $[a,b]$
e diferenciável em $(a,b)$. Então existe um ponto $c \in (a,b)$
tal que
\begin{equation*}
f(b)-f(a) = f'(c)(b-a) .
\end{equation*}
\end{thm}

Para uma interpretação geométrica do teorema do valor médio, veja a
\figureref{mvtfig}. A ideia é que o valor $\frac{f(b)-f(a)}{b-a}$
é a inclinação da reta que passa pelos pontos $\bigl(a,f(a)\bigr)$
e $\bigl(b,f(b)\bigr)$.
Então $c$ é o ponto em que $f'(c) = \frac{f(b)-f(a)}{b-a}$, isto
é, a reta tangente no ponto $\bigl(c,f(c)\bigr)$ tem a mesma inclinação que a
reta que passa por $\bigl(a,f(a)\bigr)$ e $\bigl(b,f(b)\bigr)$.
O nome vem do fato de que a inclinação da reta secante
é o valor médio da derivada; portanto, o valor médio da derivada é atingido
no interior do intervalo.

O teorema decorre do \hyperref[thm:rolle]{teorema de Rolle}
subtraindo de $f$ a função afim
que assume os mesmos valores que $f$ em $a$ e $b$.
O gráfico dessa função afim é a
reta secante — a reta que passa por
$\bigl(a,f(a)\bigr)$ e $\bigl(b,f(b)\bigr)$ — e sua derivada é
$\frac{f(b)-f(a)}{b-a}$.
Então procuramos um ponto em que essa
diferença tenha derivada nula.

\begin{myfigureht}
\myincludegraphics{mvtfig}{%
Gráfico de uma função no intervalo de a a b.
A função se parece com a anterior, mas agora não vale zero
nas extremidades: é negativa na extremidade esquerda
e positiva na direita. Uma reta que passa pelas extremidades
da curva está desenhada, isto é, uma reta que passa pelos pontos
(a, f de a) e (b, f de b). No ponto correspondente a
c no eixo horizontal, a reta tangente à função é
paralela à reta que passa pelas extremidades.}
\caption{Interpretação gráfica do teorema do valor médio.\label{mvtfig}}
\end{myfigureht}


\begin{proof}
Defina a
função $g \colon [a,b] \to \R$ por
\begin{equation*}
g(x) \coloneqq
f(x)-
\left( f(b)+\frac{f(b)-f(a)}{b-a}(x-b) \right)
=
f(x)-
f(b)-\frac{f(b)-f(a)}{b-a}(x-b) .
\end{equation*}
A função $g$ é diferenciável em $(a,b)$
e contínua em $[a,b]$, além de satisfazer $g(a) = 0$ e $g(b) = 0$. Assim, existe
$c \in (a,b)$ tal que $g'(c) = 0$, isto é,
\begin{equation*}
0 = g'(c) = f'(c)-\frac{f(b)-f(a)}{b-a} .
\end{equation*}
Em outras palavras,
$f(b)-f(a) = f'(c)(b-a)$.
\end{proof}

A demonstração pode ser generalizada. Considerando
$g(x) \coloneqq
\bigl(f(x)-f(b)\bigr)
\bigl(\varphi(b)-\varphi(a)\bigr)
-
\bigl(f(b)-f(a)\bigr)
\bigl(\varphi(x)-\varphi(b)\bigr)$,
podemos provar a seguinte versão. Deixamos a demonstração como exercício.

\begin{thm}[Teorema do valor médio de Cauchy] \label{thm:cauchymvt}
\index{teorema do valor médio de Cauchy}
Sejam $f \colon [a,b] \to \R$ e $\varphi \colon [a,b] \to \R$ funções contínuas
e diferenciáveis em $(a,b)$. Então existe um ponto $c \in (a,b)$
tal que
\begin{equation*}
\bigl(f(b)-f(a)\bigr)\varphi'(c) = f'(c)\bigl(\varphi(b)-\varphi(a)\bigr) .
\end{equation*}
\end{thm}

O teorema do valor médio tem a particularidade de ser um dos poucos teoremas
citados
em tribunais. Por exemplo, quando a polícia mede a velocidade dos carros por
aeronaves ou por câmeras que leem placas, ela
mede o tempo que o carro leva para percorrer a distância entre dois pontos.
Então, o teorema do valor médio
afirma que, em algum ponto do percurso, o carro deve ter atingido a velocidade
obtida dividindo-se a diferença entre as distâncias pela diferença entre os tempos.



\subsection{Aplicações}

Vejamos algumas aplicações do teorema do valor médio.
Elas mostram o uso típico do teorema: eliminar um limite encontrando pontos
adequados em que a derivada não é apenas próxima de algum quociente de diferenças,
mas efetivamente igual a ele.
Primeiro, resolveremos nossa primeira equação diferencial.

\begin{prop} \label{prop:derzeroconst}
Seja $I$ um intervalo e
seja $f \colon I \to \R$ uma função diferenciável tal que $f'(x) = 0$
para todo $x \in I$.
Então $f$ é constante.
\end{prop}

\begin{proof}
Tomemos $x,y \in I$ arbitrários, com $x < y$.
Como $I$ é um intervalo, $[x,y] \subset I$.
A restrição de $f$ a $[x,y]$ satisfaz as hipóteses
do \hyperref[thm:mvt]{teorema do valor médio}.
Portanto, existe $c \in (x,y)$ tal que
\begin{equation*}
f(y)-f(x) = f'(c)(y-x).
\end{equation*}
Como $f'(c) = 0$, temos $f(y) = f(x)$. Logo,
a função é constante.
\end{proof}

Agora que sabemos o que significa uma função permanecer constante, vamos
considerar as funções crescentes e decrescentes.
Dizemos que $f \colon I \to \R$ é \emph{\myindex{crescente}}
(respectivamente, \emph{\myindex{estritamente crescente}}) se
$x < y$ implica $f(x) \leq f(y)$ (respectivamente, $f(x) < f(y)$).
Definimos
\emph{\myindex{decrescente}} e
\emph{\myindex{estritamente decrescente}} da mesma maneira, invertendo as
desigualdades para $f$.

\begin{prop} \label{incdecdiffprop}
Seja $I$ um intervalo e
seja $f \colon I \to \R$ uma função diferenciável.
%\begin{enumerate}[(i),itemsep=0.5\itemsep,parsep=0.5\parsep,topsep=0.5\topsep,partopsep=0.5\partopsep]
\begin{enumerate}[(i)]
\item $f$ é crescente se, e somente se, $f'(x) \geq 0$ para todo $x \in I$.
\item $f$ é decrescente se, e somente se, $f'(x) \leq 0$ para todo $x \in I$.
\end{enumerate}
\end{prop}

\begin{proof}
Demonstremos o primeiro item. Suponha que $f$ seja crescente.
Para todos $x,c \in I$ com $x \neq c$,
\begin{equation*}
\frac{f(x)-f(c)}{x-c} \geq 0 .
\end{equation*}
Tomando o limite quando $x$ tende a $c$, vemos que $f'(c) \geq 0$.

Para provar a outra implicação, suponha que $f'(x) \geq 0$ para todo $x \in I$.
Sejam $x, y \in I$ quaisquer, com $x < y$; note que $[x,y] \subset I$.
Pelo \hyperref[thm:mvt]{teorema do valor médio}, existe $c \in (x,y)$ tal que
\begin{equation*}
f(y)-f(x) = f'(c)(y-x) .
\end{equation*}
Como $f'(c) \geq 0$ e $y-x > 0$, temos $f(y) - f(x) \geq 0$, ou seja,
$f(x) \leq f(y)$; portanto,
$f$ é crescente.

Deixamos como exercício para o leitor o segundo item, relativo à função $f$
decrescente.
\end{proof}

Uma afirmação semelhante, mas mais fraca, vale para funções estritamente
crescentes e decrescentes.

\begin{prop} \label{incdecdiffstrictprop}
Seja $I$ um intervalo e
seja $f \colon I \to \R$ uma função diferenciável.
\begin{enumerate}[(i)]
\item
\label{incdecdiffstrictprop:i}
Se $f'(x) > 0$ para todo $x \in I$, então
$f$ é estritamente crescente.
\item
\label{incdecdiffstrictprop:ii}
Se $f'(x) < 0$ para todo $x \in I$,
então $f$ é estritamente decrescente.
\end{enumerate}
\end{prop}

A demonstração de
\ref{incdecdiffstrictprop:i}
fica como exercício.
O item \ref{incdecdiffstrictprop:ii}
decorre de
\ref{incdecdiffstrictprop:i} considerando-se $-f$.
A recíproca desta proposição não é verdadeira. A função
$f(x) \coloneqq x^3$ é estritamente crescente, mas $f'(0) = 0$.

\medskip

Outra aplicação do \hyperref[thm:mvt]{teorema do valor médio} é o seguinte
resultado sobre a localização de extremos, às vezes chamado de
\emph{\myindex{teste da primeira derivada}}.
O resultado é enunciado para um mínimo e um máximo absolutos.
Para aplicá-lo na busca de mínimos e máximos relativos,
restrinja $f$ a um intervalo $(c-\delta,c+\delta)$.

\begin{prop} \label{firstderminmaxtest}
Seja $f \colon (a,b) \to \R$ contínua. Seja $c \in (a,b)$
e suponha que
$f$ seja diferenciável em $(a,c)$ e $(c,b)$.
\begin{enumerate}[(i)]
\item Se $f'(x) \leq 0$ sempre que $x \in (a,c)$ e
 $f'(x) \geq 0$ sempre que $x \in (c,b)$, então $f$ tem um mínimo absoluto
em $c$.
\item Se $f'(x) \geq 0$ sempre que $x \in (a,c)$ e
 $f'(x) \leq 0$ sempre que $x \in (c,b)$, então $f$ tem um máximo absoluto
em $c$.
\end{enumerate}
\end{prop}

\begin{proof}
Demonstremos o primeiro item e deixemos o segundo para o leitor.
Tome $x \in (a,c)$
e uma sequência $\{ y_n\}_{n=1}^\infty$ tal que $x < y_n < c$ para todo $n$
e $\lim_{n\to\infty} y_n = c$.
Pela proposição anterior,
$f$ é decrescente em $(a,c)$, logo $f(x) \geq f(y_n)$ para todo $n$.
Como $f$ é
contínua em $c$, passando ao limite obtemos
$f(x) \geq f(c)$.

De modo semelhante, tome $x \in (c,b)$
e uma sequência $\{ y_n\}_{n=1}^\infty$ tal que $c < y_n < x$ e
$\lim_{n\to\infty} y_n = c$.
A função é crescente em $(c,b)$, logo $f(x) \geq f(y_n)$ para todo $n$.
Pela continuidade de $f$, obtemos
$f(x) \geq f(c)$. Assim, $f(x) \geq f(c)$ para todo
$x \in (a,b)$.
\end{proof}

A recíproca da proposição não vale. Veja abaixo o
\exampleref{baddifffunc:example}.

\medskip

Outra aplicação frequente do teorema do valor médio, que talvez você tenha
visto no cálculo, é o resultado a seguir sobre a diferenciabilidade nas
extremidades de um intervalo. A demonstração está no
\exerciseref{exercise:endpointderivative}.

\begin{prop} \label{prop:endpointderivative}
\leavevmode
\begin{enumerate}[(i)]
\item
Suponha que $f \colon [a,b) \to \R$ seja contínua, diferenciável em $(a,b)$
e que $\lim_{x \to a} f'(x) = L$. Então $f$ é diferenciável em $a$ e
$f'(a) = L$.
\item
Suponha que $f \colon (a,b] \to \R$ seja contínua, diferenciável em $(a,b)$
e que $\lim_{x \to b} f'(x) = L$. Então $f$ é diferenciável em $b$ e
$f'(b) = L$.
\end{enumerate}
\end{prop}

De fato, usando o resultado de extensão \propref{context:prop}, não é
necessário supor que $f$ esteja definida na extremidade. Veja o
\exerciseref{exercise:extendboundedder}.

\subsection{Continuidade das derivadas e o teorema do valor intermediário}

As derivadas de funções satisfazem a propriedade do valor intermediário.

\begin{thm}[Darboux] \label{thm:darboux} \index{teorema de Darboux}
Seja $f \colon [a,b] \to \R$ diferenciável. Suponha que $y \in \R$ seja tal que $f'(a) < y < f'(b)$ ou $f'(a) > y > f'(b)$. Então existe $c \in (a,b)$ tal que $f'(c) = y$.
\end{thm}

A demonstração consiste em definir $g$ como a diferença entre $f$ e uma função linear cuja derivada é $y$. A nova função $g$ reduz o problema ao caso $y=0$, em que $g'(a) > 0 > g'(b)$. Isto é, $g$ cresce em $a$ e decresce em $b$, portanto deve atingir um máximo no interior de $(a,b)$, onde sua derivada é nula. Veja a \figureref{darbouxthmfig}.

\begin{myfigureht}
\myincludegraphics{darbouxthmfig}{%
Gráfico de uma função com três pontos, a, c e b, nessa ordem, marcados no eixo horizontal. Retas tangentes são traçadas nos pontos correspondentes do gráfico. Da esquerda para a direita, a primeira reta tem inclinação positiva e está identificada por g linha de a maior que zero. A segunda é horizontal e está identificada por g linha de c igual a zero. A terceira tem inclinação negativa e está identificada por g linha de b menor que zero.}
\caption{Ideia da demonstração do teorema de Darboux.\label{darbouxthmfig}}
\end{myfigureht}

\begin{proof}
Suponha que $f'(a) < y < f'(b)$. Defina
\begin{equation*}
g(x) \coloneqq yx - f(x) .
\end{equation*}
A função $g$ é contínua em $[a,b]$ e, portanto, atinge um máximo em algum $c \in [a,b]$.

A função $g$ também é diferenciável em $[a,b]$. Calculando, temos $g'(x) = y-f'(x)$. Logo, $g'(a) > 0$. Como a derivada é o limite dos quocientes incrementais e é positiva, algum quociente incremental deve ser positivo. Isto é, deve existir $x > a$ tal que
\begin{equation*}
\frac{g(x)-g(a)}{x-a} > 0 ,
\end{equation*}
ou seja, $g(x) > g(a)$. Portanto, $g$ não pode atingir um máximo em $a$. Analogamente, como $g'(b) < 0$, encontramos um outro $x < b$ tal que $\frac{g(x)-g(b)}{x-b} < 0$, ou seja, $g(x) > g(b)$; assim, $g$ não pode atingir um máximo em $b$. Logo, $c \in (a,b)$ e podemos aplicar o \lemmaref{relminmax:lemma}: como $g$ atinge um máximo em $c$, temos $g'(c)=0$ e, portanto, $f'(c)=y$.

De modo semelhante, se $f'(a) > y > f'(b)$, considere $g(x) \coloneqq f(x)- yx$.
\end{proof}

Já vimos que existem funções descontínuas que possuem a propriedade do valor intermediário. Embora isso possa parecer surpreendente à primeira vista, também existem funções diferenciáveis em todos os pontos cuja derivada não é contínua.

\begin{example} \label{baddifffunc:example}
Seja $f \colon \R \to \R$ a função definida por
\begin{equation*}
f(x) \coloneqq
\begin{cases}
{\bigl( x \sin(\nicefrac{1}{x}) \bigr)}^2 & \text{if } x \neq 0, \\
0 & \text{if } x = 0.
\end{cases}
\end{equation*}
Afirmamos que $f$ é diferenciável em todos os pontos, mas $f' \colon \R \to \R$ não é contínua na origem. Além disso, $f$ tem um mínimo em $0$, mas sua derivada muda de sinal infinitas vezes nas proximidades da origem. Veja a \figureref{fig:nonc1diff}.
\begin{myfigureht}
\myincludepdft{nonc1diff_full}{%
À esquerda está o gráfico de uma função que oscila entre o eixo das abscissas e a função x ao quadrado, cujo gráfico está tracejado. À direita, está uma função que oscila cada vez mais rapidamente à medida que se aproxima da origem, lembrando o gráfico do seno de um sobre x, mas com o gráfico inteiro um pouco inclinado, de modo que também cresce.}
\caption{Uma função com derivada descontínua. A função $f$ está à esquerda e $f'$ à direita. Observe que $f(x) \leq x^2$ no gráfico da esquerda.\label{fig:nonc1diff}}
\end{myfigureht}

Demonstração: Pela definição, é imediato que $f$ tem um mínimo absoluto em $0$; sabemos que $f(x) \geq 0$ para todo $x$ e que $f(0)=0$.

Para $x \neq 0$, $f$ é diferenciável, e sua derivada é $2 \sin (\nicefrac{1}{x}) \bigl( x \sin (\nicefrac{1}{x}) - \cos(\nicefrac{1}{x}) \bigr)$. Como exercício, mostre que, para $x_n = \frac{4}{(8n+1)\pi}$, temos $\lim_{n\to\infty} f'(x_n) = -1$, e que, para $y_n = \frac{4}{(8n+3)\pi}$, temos $\lim_{n\to\infty} f'(y_n) = 1$. Portanto, $f'$ não pode ser contínua em $0$, qualquer que seja o valor de $f'(0)$.

Mostremos que $f$ é diferenciável em $0$ e que $f'(0)=0$.
Para $x \neq 0$,
\begin{equation*}
\abs{\frac{f(x)-f(0)}{x-0} - 0}
=
\abs{\frac{x^2 \sin^2(\nicefrac{1}{x})}{x}}
=
\babs{x \sin^2(\nicefrac{1}{x})}
\leq
\sabs{x} .
\end{equation*}
Evidentemente, quando $x$ tende a zero, $\sabs{x}$ tende a zero e, portanto, $\abs{\frac{f(x)-f(0)}{x-0} - 0}$ tende a zero. Logo, $f$ é diferenciável em $0$ e sua derivada nesse ponto é $0$. Um ponto essencial no cálculo acima é que $\babs{f(x)} \leq x^2$; veja também os Exercícios \ref{exercise:bndmuldiff} e \ref{exercise:diffsqueeze}.
\end{example}

Às vezes, é útil supor que a derivada de uma função diferenciável é contínua. Se $f \colon I \to \R$ é diferenciável e sua derivada $f'$ é contínua em $I$, dizemos que $f$ é \emph{\myindex{continuamente diferenciável}}\index{diferenciável!continuamente}. É comum denotar por $C^1(I)$ o conjunto das funções continuamente diferenciáveis definidas em $I$.

\subsection{Exercícios}

\begin{exercise}
Conclua a demonstração de \propref{incdecdiffprop}.
\end{exercise}

\begin{exercise}
Conclua a demonstração de \propref{firstderminmaxtest}.
\end{exercise}

\begin{exercise} \label{exercise:boundeddermeanslip}
Suponha que $f \colon \R \to \R$ seja uma função diferenciável e que $f'$ seja limitada. Demonstre que $f$ é Lipschitz contínua.
\end{exercise}

\begin{exercise}
\pagebreak[1]
Suponha que $f \colon [a,b] \to \R$ seja diferenciável e que $c \in [a,b]$. Mostre que existe uma sequência $\{ x_n \}_{n=1}^\infty$ que converge para $c$, com $x_n \neq c$ para todo $n$, tal que
\begin{equation*}
f'(c) = \lim_{n\to \infty} f'(x_n).
\avoidbreak
\end{equation*}
Observe que isso \emph{não} implica que $f'$ seja contínua (por quê?).
\end{exercise}

\begin{exercise}
Suponha que $f \colon \R \to \R$ seja uma função tal que $\babs{f(x)-f(y)} \leq \sabs{x-y}^2$ para todos $x$ e $y$. Mostre que $f(x)=C$ para alguma constante $C$. Dica: mostre que $f$ é diferenciável em todos os pontos e calcule a derivada.
\end{exercise}

\begin{exercise} \label{exercise:posderincr}
Complete a demonstração de \propref{incdecdiffstrictprop}. Isto é, suponha que $I$ seja um intervalo e que $f \colon I \to \R$ seja diferenciável, com $f'(x)>0$ para todo $x \in I$. Mostre que $f$ é estritamente crescente.
\end{exercise}

\begin{exercise}
Suponha que $f \colon (a,b) \to \R$ seja diferenciável e que $f'(x) \neq 0$ para todo $x \in (a,b)$. Suponha ainda que exista $c \in (a,b)$ tal que $f'(c)>0$. Demonstre que $f'(x)>0$ para todo $x \in (a,b)$.
\end{exercise}

\begin{exercise} \label{exercise:samediffconst}
Suponha que $f \colon (a,b) \to \R$ e $g \colon (a,b) \to \R$ sejam diferenciáveis e que $f'(x)=g\'(x)$ para todo $x \in (a,b)$. Mostre que existe uma constante $C$ tal que $f(x)=g(x)+C$.
\end{exercise}

\begin{exercise}
Demonstre a seguinte versão da \myindex{regra de L'Hôpital}\index{regra de L'Hospital}.
Suponha que $f \colon (a,b) \to \R$ e $g \colon (a,b) \to \R$ sejam diferenciáveis e que $c \in (a,b)$. Suponha que $f(c)=0$, $g(c)=0$, $g'(x) \neq 0$ para $x \neq c$ e que exista o limite de $\nicefrac{f'(x)}{g'(x)}$ quando $x$ tende a $c$. Mostre que
\begin{equation*}
\lim_{x \to c} \frac{f(x)}{g(x)} = 
\lim_{x \to c} \frac{f'(x)}{g'(x)} .
\end{equation*}
Compare com \exerciseref{exercise:simpleLHopital}.
Observação: antes de fazer qualquer outra coisa, demonstre que $g(x) \neq 0$ quando $x \neq c$.
\end{exercise}

\begin{exercise}
Seja $f \colon (a,b) \to \R$ uma função diferenciável não limitada. Mostre que $f' \colon (a,b) \to \R$ não é limitada.
Observação: é importante que $(a,b)$ seja um intervalo limitado.
\end{exercise}

\begin{exercise}
Demonstre o teorema que Rolle de fato provou em 1691:
\emph{Se $f$ é um polinômio, $f'(a)=f'(b)=0$ para certos $a<b$ e não existe $c \in (a,b)$ tal que $f'(c)=0$, então há no máximo uma raiz de $f$ em $(a,b)$, isto é, no máximo um $x \in (a,b)$ tal que $f(x)=0$.}
Em outras palavras, entre duas raízes consecutivas de $f'$ há no máximo uma raiz de $f$.
Dica: suponha que haja duas raízes e veja o que acontece.
\end{exercise}

\begin{exercise}
Suponha que $a,b \in \R$ e que $f \colon \R \to \R$ seja diferenciável, com $f'(x)=a$ para todo $x$ e $f(0)=b$. Determine $f$ e demonstre que ela é a única função diferenciável com essa propriedade.
\end{exercise}

\begin{exercise} \label{exercise:endpointderivative}
\leavevmode
\begin{enumerate}[a)]
\item
Demonstre \propref{prop:endpointderivative}.
\item 
Suponha que $f \colon (a,b) \to \R$ seja contínua e diferenciável em todos os pontos, exceto em $c \in (a,b)$, e que $\\lim_{x \to c} f'(x)=L$. Demonstre que $f$ é diferenciável em $c$ e que $f'(c)=L$.
\end{enumerate}
\end{exercise}

\begin{exercise} \label{exercise:extendboundedder}
Suponha que $f \colon (0,1) \to \R$ seja diferenciável e que $f'$ seja limitada.
\begin{enumerate}[a)]
\item
Mostre que existe uma função contínua $g \colon [0,1) \to \R$ tal que $f(x)=g(x)$ para todo $x \neq 0$.\\\
Dica: use \propref{context:prop} e \exerciseref{exercise:boundeddermeanslip}.
\item
Encontre um exemplo em que $g$ não seja diferenciável em $x=0$.
\\
Dica: considere uma função baseada em $\sin(\ln x)$ e suponha que você conheça as propriedades básicas de $\sin$ e $\ln$ do cálculo.
\item
Em vez de supor que $f'$ seja limitada, suponha que $\lim_{x \to 0} f'(x)=L$. Demonstre que $g$ existe e, além disso, é diferenciável em $0$, com $g'(0)=L$.
\end{enumerate}
\end{exercise}

\begin{exercise}
Demonstre \thmref{thm:cauchymvt}.
\end{exercise}

%%%%%%%%%%%%%%%%%%%%%%%%%%%%%%%%%%%%%%%%%%%%%%%%%%%%%%%%%%%%%%%%%%%%%%%%%%%%%%

\sectionnewpage
\section{Taylor's theorem}
\label{sec:taylor}

%mbxINTROSUBSECTION

\sectionnotes{less than a lecture (optional section)}

\subsection{Derivatives of higher orders}

When $f \colon I \to \R$ is differentiable, we obtain a function
$f' \colon I \to \R$.  The function
$f'$ is called the \emph{\myindex{first derivative}} of $f$.
If $f'$ is differentiable, we denote by
$f'' \colon I \to \R$ the derivative of $f'$.  The function $f''$
is called the \emph{\myindex{second derivative}} of $f$.
\glsadd{not:secondthirdfourthder}
We similarly obtain
$f'''$, $f''''$, and so on.
With a larger number of derivatives
the notation would get out of hand; we denote
by $f^{(n)}$ the
\glsadd{not:nthder}
\emph{$n$th derivative}\index{nth derivative@$n$th derivative} of $f$.
%
When $f$ possesses $n$ derivatives, we say $f$ is
\emph{$n$ times differentiable}\index{n times differentiable@$n$ times
differentiable}\index{differentiable!n times@$n$ times}.

\subsection{Taylor's theorem}

Taylor's theorem%
\footnote{Named for the English mathematician
\href{https://en.wikipedia.org/wiki/Brook_Taylor}{Brook Taylor}
(1685--1731).
It was found by
the Scottish mathematician
\href{https://en.wikipedia.org/wiki/James_Gregory_(mathematician)}{James Gregory}
(1638--1675).  The statement we give
is due to
\href{https://en.wikipedia.org/wiki/Lagrange}{Joseph-Louis Lagrange}
(1736--1813).} (at least the version we give here)
is a generalization of the \hyperref[thm:mvt]{mean value theorem}.
Mean value theorem says that up to a small error $f(x)$ for $x$ near $x_0$ can be
approximated by $f(x_0)$:
\begin{equation*}
f(x) = f(x_0) + f'(c)(x-x_0),
\end{equation*}
where the \myquote{error} $f'(c)(x-x_0)$
is measured in terms of the first derivative
at some point $c$ between $x$ and $x_0$.
Taylor's theorem generalizes this result to higher derivatives.
It tells us that up to a small error, any $n$
times differentiable function can be approximated at a point $x_0$
by a polynomial of degree $n$.  The
error of this approximation goes to zero faster than ${(x-x_0)}^{n}$ as
$x$ goes to $x_0$.
To see why this is a good approximation, notice that for a big $n$, 
${(x-x_0)}^n$ is very small in a small interval around $x_0$.
We will require one more (so $n+1$) derivative to write the error
as we do in the mean value theorem.

\begin{defn}
For an $n$ times differentiable function $f$ defined near a point $x_0 \in \R$, define the
$n$th order \emph{\myindex{Taylor polynomial}}%
\index{nth order Taylor polynomial@$n$th order Taylor polynomial}
for $f$ at $x_0$ as
\begin{equation*}
\begin{split}
P_n^{x_0}(x)
& \coloneqq
\sum_{k=0}^n
\frac{f^{(k)}(x_0)}{k!}{(x-x_0)}^k
\\
& =
f(x_0)
+ f'(x_0)(x-x_0)
+ \frac{f''(x_0)}{2}{(x-x_0)}^2
+ \frac{f^{(3)}(x_0)}{6}{(x-x_0)}^3
\\
& \qquad
+ \cdots
+ \frac{f^{(n)}(x_0)}{n!}{(x-x_0)}^n .
\end{split}
\end{equation*}
\end{defn}

See \figureref{fig:taylorsin} for
the odd-degree Taylor polynomials for the sine function at $x_0=0$.
The even-degree terms are all zero, as even derivatives 
of sine are again sines, which are zero at a origem.
\begin{myfigureht}
\myincludegraphics{taylorsin}{%
The graph of the sine function is given.  Also several
approximations are given.  First a straight line that is
tangent at the origin is labeled as y equals P sub 1 super 0 of x.
Second a graph of a cubic that is tangent to the graph of sine
at the origin and somewhat approximates it nearby is labeled
as y equals P sub 3 super 0 of x.
Similarly with degree 5 and degree 7 approximations.  The higher the
degree the further away from the origin is the graph reasonably close 
to the sine function, for example the degree 7 approximation seems
relatively close between two points where the sine
is zero around the origin, that is, plus and minus pi.}
\caption{The odd degree Taylor polynomials for the sine
function.\label{fig:taylorsin}}
\end{myfigureht}

The statement of Taylor's theorem that we give includes the
\hyperref[thm:mvt]{mean value theorem}, which can be thought of as
Taylor's theorem for $n=0$.

\begin{thm}[Taylor]\index{Taylor's theorem} \label{thm:taylor}
Suppose $f \colon [a,b] \to \R$ is a function with $n$ continuous
derivatives on $[a,b]$ and such that $f^{(n+1)}$ exists on $(a,b)$.
Given distinct points $x_0$ and $x$ in $[a,b]$,
there is a point $c$ between $x_0$
and $x$ such that
\begin{equation*}
f(x)=P_{n}^{x_0}(x)+\frac{f^{(n+1)}(c)}{(n+1)!}{(x-x_0)}^{n+1} .
\end{equation*}
\end{thm}

The term $R_n^{x_0}(x)\coloneqq
f(x)-P_n^{x_0}(x) =
\frac{f^{(n+1)}(c)}{(n+1)!}{(x-x_0)}^{n+1}$ is called the
\emph{remainder term}\index{remainder term in Taylor's formula}.  This
form is called the
\emph{\myindex{Lagrange form}} of the remainder.  There are other ways
to write the remainder term, but we skip those.  Note that $c$ depends on
both $x$ and $x_0$.

\begin{proof}
Let $M_{x,x_0}$ be the number (depending on $x$ and $x_0$) solving the equation
\begin{equation*}
f(x)=P_{n}^{x_0}(x)+M_{x,x_0}{(x-x_0)}^{n+1} .
\end{equation*}
Define a function $g(s)$ by
\begin{equation*}
g(s) \coloneqq f(s)-P_n^{x_0}(s)-M_{x,x_0}{(s-x_0)}^{n+1} .
\end{equation*}
We compute
the $k$th derivative at $x_0$ of the Taylor polynomial
${(P_n^{x_0})}^{(k)}(x_0) = f^{(k)}(x_0)$ for
$k=0,1,2,\ldots,n$ (the zeroth derivative of a function is the function
itself).  Therefore,
\begin{equation*}
g(x_0) = g'(x_0) = g''(x_0) = \cdots = g^{(n)}(x_0) = 0 .
\end{equation*}
In particular, $g(x_0) = 0$.
We also have $g(x) = 0$.
By the
\hyperref[thm:mvt]{mean value theorem},
there exists an $x_1$ between $x_0$ and $x$ such that $g'(x_1) = 0$.
Applying the \hyperref[thm:mvt]{mean value theorem}
to $g'$, we obtain that there exists
$x_2$ between $x_0$ and $x_1$ (and therefore between $x_0$ and $x$)
such that $g''(x_2) = 0$.  We repeat the
argument $n+1$ times to obtain a number $x_{n+1}$ between $x_0$ and $x_n$
(and therefore between $x_0$ and $x$) such that $g^{(n+1)}(x_{n+1}) = 0$.

Let $c \coloneqq x_{n+1}$.
We compute the $(n+1)$th derivative of $g$ to find
\begin{equation*}
g^{(n+1)}(s) = f^{(n+1)}(s)-(n+1)!\,M_{x,x_0} .
\end{equation*}
Plugging in $c$ for $s$, we obtain $M_{x,x_0} = \frac{f^{(n+1)}(c)}{(n+1)!}$, and
we are done.
\end{proof}

In the proof, we found
${(P_n^{x_0})}^{(k)}(x_0) = f^{(k)}(x_0)$ for $k=0,1,2,\ldots,n$.
Therefore, the Taylor polynomial has the same derivatives as $f$ at $x_0$
up to the $n$th derivative.  That is why the Taylor polynomial is
a good approximation to $f$.
Notice how in \figureref{fig:taylorsin} the Taylor polynomials are
reasonably good approximations to the sine near $x=0$.

We do not necessarily get good approximations
by the Taylor polynomial everywhere.
Consider expanding the function
$f(x) \coloneqq \frac{x}{1-x}$ around 0,
for $x < 1$, we get the graphs in
\figureref{fig:taylorgeom}.  The dotted lines are the first, second, and
third degree approximations.
The dashed line is the 20th degree polynomial.
The approximation does seem to get
better as the degree rises for $x > -1$.
For $x < -1$, it in fact gets visibly worse.
The polynomials
are the partial sums of the geometric series $\sum_{n=1}^\infty x^n$,
and the series only converges on $(-1,1)$.
See the discussion of power series in
\sectionref{sec:moreonseries}.

\begin{myfigureht}
\myincludegraphics{taylorgeom}{%
In bold line is a graph of a function on the interval
from minus 2 to 1, that comes in flat from the right from below
the x-axis, crosses the x-axis at the origin, and then turns sharply
upward leaving the picture before we get to 1.
The graphs of the first three Taylor polynomials are shown in
dotted line.  That is,
a straight line, a quadratic, and a cubic.
They all seem to be 
somewhat reasonable approximation very near the origin
but get worse as we get further.
A dashed line signifying the degree 20 approximation is drawn and
from some point slightly more than minus 1 until where the bold
line leaves the picture on the right, we cannot tell the difference
between it and the bold line.  However, right around minus 1
if we move from right to left, it shoots upwards leaving the picture,
so for x less than minus 1 the approximation would be ludicrously worse
than the small degree approximations (which are not good already).}
\caption{The function $\frac{x}{1-x}$, and the Taylor polynomials
$P_1^0$, $P_2^0$, $P_3^0$ (all dotted), and the polynomial $P_{20}^0$
(dashed).\label{fig:taylorgeom}}
\end{myfigureht}

If $f$ is \emph{\myindex{infinitely
differentiable}}\index{differentiable!infinitely},
that is, if $f$ can be
differentiated any number of times, then 
we define the \emph{\myindex{Taylor series}}:
\begin{equation*}
%T^{x_0}(x)
%\coloneqq
\sum_{k=0}^\infty
\frac{f^{(k)}(x_0)}{k!}{(x-x_0)}^k .
\end{equation*}
There is no guarantee that this series converges for any
$x \neq x_0$.  Even where it does converge, there is no guarantee
that it converges to the function $f$.  Functions $f$
whose Taylor series at every point $x_0$
converges to $f$ in some open interval containing $x_0$
are called
\emph{analytic functions}\index{analytic function}.
Many functions one tends to see in practice are analytic.
See \exerciseref{exercise:nonanalytic}, for an example of a non-analytic
function.

\medskip

The definition of derivative says that
a function is
differentiable if it
is locally approximated by a line.
We mention in passing that there exists a converse to Taylor's
theorem,
which we will neither state nor prove,
saying that if a function is
locally approximated in a certain way by a polynomial of degree $d$, then it
has $d$ derivatives.

\medskip

Taylor's theorem gives a quick proof of a version of
the second derivative test.
By a \emph{\myindex{strict relative minimum}}\index{minimum!strict relative}
of $f$ at $c$,
we mean that there exists a $\delta > 0$ such that $f(x) > f(c)$ for
all $x \in (c-\delta,c+\delta)$ where $x\neq c$.
A \emph{\myindex{strict relative maximum}}\index{maximum!strict relative}
is defined similarly.
Continuity of the second derivative is not needed, but the proof is more
difficult and is left as an exercise.  The proof also generalizes
into the $n$th derivative test, which is also left as
an exercise.

\begin{prop}[Second derivative test]\index{second derivative test}
Suppose $f \colon (a,b) \to \R$ is twice continuously differentiable,
$x_0 \in (a,b)$, $f'(x_0) = 0$ and $f''(x_0) > 0$.  Then $f$ has a strict relative
minimum at $x_0$.
\end{prop}

\begin{proof}
As $f''$ is continuous, there exists a $\delta > 0$
such that $f''(c) > 0$ for all $c \in (x_0-\delta,x_0+\delta)$,
see \exerciseref{exercise:positivecontneigh}.
Take $x \in (x_0-\delta,x_0+\delta)$, $x \neq x_0$.
Taylor's theorem says that for some $c$ between $x_0$ and $x$,
\begin{equation*}
f(x) 
=
f(x_0) + f'(x_0) (x-x_0) +
\frac{f''(c)}{2}{(x-x_0)}^{2} 
=
f(x_0) + \frac{f''(c)}{2}{(x-x_0)}^{2}  .
\end{equation*}
As $f''(c) > 0$ and ${(x-x_0)}^2 > 0$, we have $f(x) > f(x_0)$.
\end{proof}

\subsection{Exercises}

\begin{exercise}
Compute the $n$th Taylor polynomial at $0$ for the exponential function.
\end{exercise}

\begin{exercise}
Suppose $p$ is a polynomial of degree $d$.  Given $x_0 \in \R$,
show that
the $d$th Taylor polynomial for $p$ at $x_0$ is equal to $p$.
\end{exercise}

\begin{exercise}
Let $f(x) \coloneqq \sabs{x}^3$.  Compute $f'(x)$ and $f''(x)$ for all $x$,
but show that $f^{(3)}(0)$ does not exist.
\end{exercise}

\begin{exercise}
Suppose $f \colon \R \to \R$ has $n$ continuous derivatives.  Show
that for every $x_0 \in \R$,
there exist polynomials $P$ and $Q$ of degree $n$ and 
an $\epsilon > 0$ such that $P(x) \leq f(x) \leq Q(x)$ for all $x \in
[x_0,x_0+\epsilon]$  and
$Q(x)-P(x) = \lambda {(x-x_0)}^n$ for some $\lambda \geq 0$.
\end{exercise}

\begin{exercise}
If $f \colon [a,b] \to \R$ has $n+1$ continuous derivatives\footnote{%
The same statement holds if we only require $n$ derivatives,
but it is quite a bit harder to prove.
}
and $x_0 \in [a,b]$,
prove
$\lim\limits_{x\to x_0} \frac{R_n^{x_0}(x)}{{(x-x_0)}^n} = 0$.
\end{exercise}

\begin{exercise}
Suppose $f \colon [a,b] \to \R$ has $n+1$ continuous derivatives
and $x_0 \in (a,b)$.
Prove: $f^{(k)}(x_0) = 0$ for all $k = 0, 1, 2, \ldots, n$
if and only if $\lim\limits_{x\to x_0} \frac{f(x)}{{(x-x_0)}^{n+1}}$ exists.
\end{exercise}

\begin{exercise}
Suppose $a,b,c \in \R$ and $f \colon \R \to \R$ is twice differentiable,
$f''(x) = a$ for all $x$, $f'(0) = b$, and $f(0) = c$.  Find $f$ and prove that 
it is the unique differentiable function with this property.
\end{exercise}

\begin{exercise}[Challenging]
Show that a simple converse to Taylor's theorem does not hold.
Find a function $f \colon \R \to \R$ with no second derivative at $x=0$ such that
$\babs{f(x)} \leq \babs{x^3}$, that is, $f$ goes to zero at 0 faster than $x^2$, and
while $f'(0)$ exists, $f''(0)$ does not.
\end{exercise}

\begin{exercise} \label{exercise:extendboundedder2}
Suppose $f \colon (0,1) \to \R$ is differentiable and $f''$
is bounded.
\begin{enumerate}[a)]
\item
Show that there exists a once differentiable function $g \colon [0,1) \to \R$
such that $f(x) = g(x)$ for all $x \neq 0$.  Hint: 
See
\exerciseref{exercise:extendboundedder}.
\item
Find an example where the $g$ is not twice differentiable at $x=0$.
\end{enumerate}
\end{exercise}

\begin{exercise}
Prove the
\emph{$n$th derivative test}\index{nth derivative test@$n$th derivative test}.
Suppose $n \in \N$,
$x_0 \in (a,b)$, and $f \colon (a,b) \to \R$ is $n$ times continuously
differentiable, with $f^{(k)}(x_0) = 0$ for $k=1,2,\ldots,n-1$, and
$f^{(n)}(x_0)
\neq 0$.
Prove:
\begin{enumerate}[a)]
\item
If $n$ is odd, then $f$ has neither a relative minimum,
nor a maximum at $x_0$.
\item
If $n$ is even, then $f$ has a strict relative minimum at $x_0$ if
$f^{(n)}(x_0) > 0$
and a strict relative maximum at $x_0$ if $f^{(n)}(x_0) < 0$.
\end{enumerate}
\end{exercise}

\begin{exercise}
Prove the more general version of the second derivative test.
Suppose $f \colon (a,b) \to \R$ is differentiable and $x_0 \in (a,b)$
is such that, $f'(x_0) = 0$, $f''(x_0)$ exists, and $f''(x_0) > 0$.
Prove that $f$ has a strict relative
minimum at $x_0$.  Hint: Consider the limit definition of $f''(x_0)$.
\end{exercise}

%%%%%%%%%%%%%%%%%%%%%%%%%%%%%%%%%%%%%%%%%%%%%%%%%%%%%%%%%%%%%%%%%%%%%%%%%%%%%%

\sectionnewpage
\section{Inverse function theorem}
\label{sec:ift}

%mbxINTROSUBSECTION

\sectionnotes{less than 1 lecture (optional section, needed for
\sectionref{sec:logandexp}, requires 
\sectionref{sec:monotonefunc})}

\subsection{Inverse function theorem}

We start with a simple example.  Consider the function $f(x) \coloneqq a x$ for a
number $a \neq 0$.  Then $f \colon \R \to \R$ is bijective, and the inverse
is $f^{-1}(y) = \frac{1}{a} y$.  In particular, $f'(x) = a$ and 
$(f^{-1})'(y) = \frac{1}{a}$.  As differentiable functions are
\myquote{infinitesimally like} linear functions, we expect the same
sort of behavior from the inverse of a differentiable function.
The main idea of differentiating inverse functions is the following lemma.

\begin{lemma} \label{lemma:ift}
Let $I,J \subset \R$ be intervals.
If $f \colon I \to J$ is strictly monotone (hence one-to-one),
onto ($f(I) = J$),
differentiable at $x_0 \in I$, and $f'(x_0) \neq 0$,
then the inverse 
$f^{-1}$ is differentiable at $y_0 = f(x_0)$ and
\begin{equation*}
(f^{-1})'(y_0) = \frac{1}{f'\bigl( f^{-1}(y_0) \bigr)} = \frac{1}{f'(x_0)} .
\end{equation*}
If $f$ is continuously differentiable and $f'$ is never zero, then $f^{-1}$
is continuously differentiable.
\end{lemma}

\begin{proof}
By \propref{prop:invcont}, $f$ has a continuous inverse.  For convenience,
call the inverse $g \colon J \to I$.
Let $x_0,y_0$ be as in the statement.  For $x \in I$, write $y \coloneqq f(x)$.
If $x \neq x_0$, and so $y \neq y_0$, we find
\begin{equation*}
\frac{g(y)-g(y_0)}{y-y_0} =
\frac{g\bigl(f(x)\bigr)-g\bigl(f(x_0)\bigr)}{f(x)-f(x_0)} =
\frac{x-x_0}{f(x)-f(x_0)} .
\end{equation*}
See \figureref{inversefig} for the geometric idea.
\begin{myfigureht}
\myincludepdft{inversefigAB}{%
Two diagrams.
On the left, a function f that looks like the square root
is shown in the first quadrant.  On the horizontal axis
there is a point labeled x equals g of y.  A dashed
line goes vertically up from this point until it hits the corresponding
point on the graph of f.  Another dashed line, now horizontal goes left
until it hits the y axis and a point marked f of x equals y.
Another point x sub 0 equals g of y sub 0 is also marked
a bit to the right of the first.  The same setup with the
dashed lines is repeated except instead of x it is x sub 0 and
instead of y it is y sub 0.
A secant line is drawn through the two points on the graph
and it is labeled as having slope
f of x minus f of x sub 0, the whole thing over quantity
x minus x sub 0
and that equals
y minus y sub 0, the whole thing over quantity
g of y minus g of y sub 0.
The right hand diagram is the same thing but the picture
is simply reflected across the x equals y line, so the x-axis
becomes the y-axis and vice-versa, and the graph now looks more
like x squared.  The two expressions of the tangent
line are the reciprocals.}
\caption{Interpretation of the derivative of the inverse
function.\label{inversefig}}
\end{myfigureht}

Let
\begin{equation*}
Q(x) \coloneqq
\begin{cases}
\frac{x-x_0}{f(x)-f(x_0)} & \text{if } x \neq x_0, \\
\frac{1}{f'(x_0)}         & \text{if } x = x_0 \quad \text{(notice that }
f'(x_0) \neq 0 \text{)}.
\end{cases}
\end{equation*}
As $f$ is differentiable at $x_0$, 
\begin{equation*}
\lim_{x \to x_0} Q(x) =
\lim_{x \to x_0} 
\frac{x-x_0}{f(x)-f(x_0)} 
=
\frac{1}{f'(x_0)}
=
Q(x_0) ,
\end{equation*}
that is, $Q$ is continuous at $x_0$.
As $g(y)$ is continuous at $y_0$,
the composition $Q\bigl(g(y)\bigr) = \frac{g(y)-g(y_0)}{y-y_0}$
is continuous at $y_0$ by
\propref{prop:compositioncont}.
Therefore,
\begin{equation*}
\frac{1}{f'\bigl(g(y_0)\bigr)}
= Q\bigl(g(y_0)\bigr)
= \lim_{y \to y_0} Q\bigl(g(y)\bigr)
= \lim_{y \to y_0} \frac{g(y)-g(y_0)}{y-y_0} .
\end{equation*}
So $g$ is differentiable at $y_0$, and $g'(y_0) =
\frac{1}{f'\left(\vphantom{1^1_1}g(y_0)\right)}$.

If $f'$ is continuous and nonzero at all $x \in I$,
then the lemma applies at all $x \in I$.  As $g$ is also
continuous (it is differentiable), the derivative $g'(y) =
\frac{1}{f'\left(\vphantom{1^1_1}g(y)\right)}$ must be continuous.
\end{proof}

What is usually called the inverse function theorem is the following result.

\begin{thm}[Inverse function theorem]\index{inverse function theorem}
Let $f \colon (a,b) \to \R$ be a continuously differentiable function,
$x_0 \in (a,b)$ a point where $f'(x_0) \neq 0$.  Then there exists
an open interval $I \subset (a,b)$ with $x_0 \in I$, the
restriction $f|_{I}$ is injective with a continuously differentiable inverse
$g \colon J \to I$ defined on an interval $J \coloneqq f(I)$,
and
\begin{equation*}
g'(y) = \frac{1}{f'\bigl( g(y) \bigr)} \qquad \text{for all } y \in J.
\end{equation*}
\end{thm}

\begin{proof}
Without loss of generality, suppose $f'(x_0) > 0$.  As $f'$ is
continuous, there must exist an open interval $I = (x_0-\delta,x_0+\delta)$ 
such that $f'(x) > 0$ for all $x \in I$.  See
\exerciseref{exercise:positivecontneigh}.

By \propref{incdecdiffstrictprop}, $f$ is strictly increasing
on $I$, and hence the restriction $f|_{I}$ is bijective onto $J \coloneqq f(I)$.
As $f$ is continuous,
\corref{cor:continterval}
(or directly via the
\hyperref[IVT:thm]{intermediate value theorem})
implies that
$f(I)$ is an interval.
Now apply \lemmaref{lemma:ift}.
\end{proof}

In \exampleref{example:sqrt2}, we saw how difficult an
endeavor was proving the existence of $\sqrt{2}$ without any
tools.
With the \hyperref[IVT:thm]{intermediate value theorem},
the existence of roots is almost trivial, and
with the machinery of this section, we will prove
far more than mere existence.

\begin{cor}
Given $n \in \N$ and $x \geq 0$, there exists a unique 
number $y \geq 0$ (denoted $x^{1/n} \coloneqq y$), such that $y^n = x$.  Furthermore,
the function $g \colon (0,\infty) \to (0,\infty)$ defined by
$g(x) \coloneqq x^{1/n}$ is continuously differentiable and
\begin{equation*}
g'(x) = \frac{1}{nx^{(n-1)/n}} = \frac{1}{n} \, x^{(1-n)/n} ,
\end{equation*}
using the convention $x^{m/n} \coloneqq {(x^{1/n})}^{m}$.
\end{cor}

\begin{proof}
For $x=0$, the existence of a unique root is trivial.

Let $f \colon (0,\infty) \to (0,\infty)$ be defined by $f(y) \coloneqq y^n$.
The function $f$ is continuously differentiable,
and $f'(y) = ny^{n-1}$, see \exerciseref{exercise:diffofxn}.
For $y > 0$, the derivative $f'$ is strictly positive,
and so again by \propref{incdecdiffstrictprop}, $f$ is strictly
increasing (this can also be proved directly) and hence injective.
Suppose $M$ and $\epsilon$ are such that
$M > 1$ and $1 > \epsilon > 0$.
Then
$f(M) = M^n \geq M$ and
$f(\epsilon) = \epsilon^n \leq \epsilon$.
For every $x$ with $\epsilon < x < M$,
we have, by the
\hyperref[IVT:thm]{intermediate value theorem}, that $x \in
f\bigl( [\epsilon,M] \bigr) \subset
f\bigl( (0,\infty) \bigr)$.  As $M$ and $\epsilon$ were arbitrary, $f$ is onto
$(0,\infty)$, and hence $f$ is bijective.
Let $g$ be the inverse of $f$, and we obtain
the existence and uniqueness of positive
$n$th roots.  \lemmaref{lemma:ift} says that $g$ has a continuous
derivative and $g'(x) =
\frac{1}{f'\left(\vphantom{1^1_1}g(x)\right)} = \frac{1}{n {(x^{1/n})}^{n-1}}$.
\end{proof}

\begin{example}
The corollary provides a good example of where the inverse function theorem
gives us an interval smaller than $(a,b)$.  Take $f \colon \R \to \R$
defined by $f(x) \coloneqq x^2$.  Then $f'(x_0) \neq 0$
as long as $x_0 \neq 0$.  If $x_0 > 0$, we can take $I=(0,\infty)$, but
no larger.
\end{example}

\begin{example}
Another useful example is $f(x) \coloneqq x^3$.  The function $f \colon \R \to \R$ is
one-to-one and onto, so $f^{-1}(y) = y^{1/3}$ exists on the entire real
line, including zero and negative $y$.  The function $f$ has
a continuous derivative, but $f^{-1}$ has no derivative at a origem.  The
point is that $f'(0) = 0$.  See \figureref{cubecuberootfig} for a graph.
Notice the vertical tangent on the cube root at a origem.
See also \exerciseref{exercise:oddroot}.
\begin{myfigureht}
\myincludegraphics{cubecuberoot}{%
The graph y equals the cube root of x in bold line
and the graph of y equals x cubed in dashed line.
We note that while the graph of x cubed is tangent to the x-axis
at the origin, the cube root is a reflection across the x equals y
line and is therefore tangent to the vertical y-axis at a origem.}
\caption{Graphs of $x^3$ and $x^{1/3}$.\label{cubecuberootfig}}
\end{myfigureht}
\end{example}


\subsection{Exercises}

\begin{exercise}
Suppose $f \colon \R \to \R$ is continuously differentiable and
$f'(x) > 0$ for all $x$.  Show that $f$ is invertible on the interval $J =
f(\R)$, the inverse is continuously differentiable, and ${(f^{-1})}'(y) >
0$ for all $y \in f(\R)$.
\end{exercise}

\begin{exercise}
\pagebreak[2]
Suppose $I,J$ are intervals and a monotone onto $f \colon I \to J$ has an inverse $g \colon J \to I$.
Suppose you already know that both $f$ and $g$ are differentiable
everywhere and $f'$ is never zero.  Using the chain rule but not
\lemmaref{lemma:ift}, prove the
formula $g'(y) = \frac{1}{f'\left(\vphantom{1_1^1}g(y)\right)}$. % \bigl( \bigr) are overly big here!
Remark: This exercise is the same as \exerciseref{exercise:inversederformula},
no need to do it again if you have
solved it already.
\end{exercise}

\begin{exercise}
\pagebreak[2]
Let $n\in \N$ be even.
Prove that every $x > 0$ has a unique negative $n$th root.
That is, there exists a negative number $y$ such that $y^n = x$.
Compute the derivative
of the function $g(x) \coloneqq y$.
\end{exercise}

\begin{exercise} \label{exercise:oddroot}
Let $n \in \N$ be odd and $n \geq 3$.
Prove that every $x$ has a unique $n$th root.
That is, there exists a number $y$ such that $y^n = x$.  Prove that
the function defined by $g(x) \coloneqq y$ is differentiable except at $x=0$
and compute the derivative.  Prove that $g$ is not differentiable at $x=0$.
\end{exercise}

\begin{exercise}[requires \sectionref{sec:taylor}]
Show that if in the inverse function theorem $f$ has $k$ continuous
derivatives, then the inverse function $g$ also has $k$ continuous
derivatives.
\end{exercise}

\begin{exercise}
Let $f(x) \coloneqq x + 2 x^2 \sin(\nicefrac{1}{x})$ for $x \neq 0$ and
$f(0) \coloneqq 0$.  Show that $f$ is differentiable at all $x$, that $f'(0) > 0$,
but that $f$ is not invertible
on any open interval containing a origem.
\end{exercise}

\begin{exercise}
\leavevmode
\begin{enumerate}[a)]
\item
Let $f \colon \R \to \R$ be a continuously differentiable function
and $k > 0$ be a number such that $f'(x) \geq k$ for all $x \in \R$.
Show $f$ is one-to-one and onto, and has a continuously differentiable
inverse $f^{-1} \colon \R \to \R$.
\item
Find an example $f \colon \R \to \R$
where $f'(x) > 0$
for all $x$, but $f$ is not onto.
\end{enumerate}
\end{exercise}

\begin{exercise}
Suppose $I,J$ are intervals and a monotone onto $f \colon I \to J$ has an inverse $g \colon J \to I$.
Suppose $x \in I$ and $y \coloneqq f(x) \in J$, and that $g$ is differentiable at
$y$.  Prove:
\begin{enumerate}[a)]
\item
If $g'(y) \neq 0$, then $f$ is differentiable at $x$.
\item
If $g'(y) = 0$, then $f$ is not differentiable at $x$.
\end{enumerate}
\end{exercise}
